\documentclass[11pt,a4paper]{article}

\usepackage[utf8]{inputenc}
\usepackage[T1]{fontenc}
\usepackage[english]{babel}
\usepackage{amsmath,amssymb,amsthm}
\usepackage{geometry}
\usepackage{hyperref}
\usepackage{booktabs}
\usepackage{array}
\usepackage{tabularx}
\usepackage{xcolor}
\usepackage{mdframed}
\usepackage{enumitem}
\usepackage{microtype}
\usepackage{graphicx}
\usepackage{authblk}
\usepackage{fancyhdr}
\usepackage{titlesec}
\usepackage[numbers,sort&compress]{natbib}

\geometry{margin=2.5cm}

\hypersetup{
  colorlinks=true,
  linkcolor=blue!60!black,
  citecolor=blue!60!black,
  urlcolor=blue!60!black,
  pdftitle={The Epistemic Interface Problem in Multi-Agent Systems},
  pdfauthor={Andrei Dinga-Reassi},
}

%% Theorem environments
\theoremstyle{plain}
\newtheorem{theorem}{Theorem}[section]
\newtheorem{lemma}[theorem]{Lemma}
\newtheorem{corollary}[theorem]{Corollary}
\newtheorem{proposition}[theorem]{Proposition}

\theoremstyle{definition}
\newtheorem{definition}[theorem]{Definition}

\theoremstyle{remark}
\newtheorem{remark}[theorem]{Remark}

%% Framed box for key results
\newmdenv[
  linecolor=blue!40!black,
  backgroundcolor=blue!3,
  linewidth=1pt,
  innertopmargin=8pt,
  innerbottommargin=8pt,
  innerleftmargin=10pt,
  innerrightmargin=10pt,
]{keyresult}

\newmdenv[
  linecolor=gray!60,
  backgroundcolor=gray!5,
  linewidth=0.5pt,
  innertopmargin=6pt,
  innerbottommargin=6pt,
  innerleftmargin=8pt,
  innerrightmargin=8pt,
]{changebox}

%% Header
\pagestyle{fancy}
\fancyhf{}
\rhead{\small EIP — Working Paper v0.2 — May 2026}
\lhead{\small Dinga-Reassi}
\cfoot{\thepage}

%% Section formatting
\titleformat{\section}{\large\bfseries}{\thesection.}{0.5em}{}
\titleformat{\subsection}{\normalsize\bfseries}{\thesubsection.}{0.5em}{}

% Math shortcuts
\newcommand{\Real}{\mathbb{R}}
\newcommand{\Bel}{\mathbb{B}}
\newcommand{\Theta}{\Theta}
\newcommand{\norm}[1]{\left\|#1\right\|}

\begin{document}

%% ─────────────────────────────────────────────
%%  TITLE PAGE
%% ─────────────────────────────────────────────
\begin{titlepage}
\centering
\vspace*{1.5cm}

{\LARGE\bfseries The Epistemic Interface Problem\\[0.4em]
in Multi-Agent Systems}\\[0.6em]
{\large\itshape Why Latent-Space Communication Requires a Formal Contract}

\vspace{0.8cm}
{\large Version 0.2 --- Working Paper --- May 2026}

\vspace{0.6cm}
\begin{tabular}{c}
\textbf{Andrei Dinga-Reassi}\\[0.2em]
Independent Researcher\\[0.2em]
\href{https://github.com/dravitch}{github.com/dravitch}\\[0.2em]
\texttt{}
\end{tabular}

\vspace{0.6cm}
{\small
Zenodo: \href{https://doi.org/10.5281/zenodo.20478868}{10.5281/zenodo.20478868}
\quad$\cdot$\quad
Repository: \href{https://github.com/dravitch/eip-theorem}{github.com/dravitch/eip-theorem}
}

\vspace{0.8cm}
\hrule
\vspace{0.4cm}
\begin{minipage}{0.9\textwidth}
\small\itshape
Based on: RecursiveMAS \cite{yang2026} $\cdot$ IB \cite{tishby1999} $\cdot$
GIB \cite{westphal2025} $\cdot$ Deep EK-NN \cite{hoarau2025} $\cdot$
Epistemic DL \cite{manchingal2022} $\cdot$ IoA \cite{ioa2025} $\cdot$
Perry \& Tsoukias \cite{perry1998} $\cdot$ Smets \cite{smets1994} $\cdot$
Belnap \cite{belnap1977} $\cdot$ Deno\oe{}ux \cite{denoeux2008}
\end{minipage}
\vspace{0.4cm}
\hrule

\vspace{0.8cm}
\begin{changebox}
\small
\textbf{Key changes v0.1 $\to$ v0.2.}
Section~3.3 corrected: the IB connection is reformulated as a neural-implementation
phenomenon (not predicted by Tishby's original IB).
Section~5.3 strengthened: complete axiomatic justification of rules O1/O3/O5
(LCP, Klawonn \& Smets 1992, Deno\oe{}ux 2008).
Section~5.5 added: theoretical unification of \texttt{belief\_mass} and
\texttt{belnap\_state} via Perry \& Tsoukias's (1998) fuzzy Belnap space.
Section~10 added: future extensions including latent credal channel, adaptive
recursive fusion, and epistemic-measure stopping criterion.
\end{changebox}

\vspace{0.6cm}
\begin{mdframed}[linecolor=gray!40,backgroundcolor=gray!4]
\small\textbf{Acknowledgements.}
This work was developed with the assistance of Claude (Anthropic) as a
collaborative reasoning tool for specification design, formal verification, and
peer review simulation. The intellectual contributions, design decisions, and
responsibility for the content are those of the author.
Experimental validation was conducted in the companion repository
\href{https://github.com/symbioticode/epistemic-impossibility-validation}%
{symbioticode/epistemic-impossibility-validation}.
\end{mdframed}
\end{titlepage}

%% ─────────────────────────────────────────────
%%  ABSTRACT -- page 2
%% ─────────────────────────────────────────────
\begin{abstract}
We identify and formalise a fundamental gap in the multi-agent systems (MAS)
literature, which we call the \textbf{Epistemic Interface Problem (EIP)}: the
absence of a formal contract governing the transformation of an agent's internal
latent state into an auditable epistemic declaration. Current research conflates
two distinct problems --- \emph{expressivity} (how to transmit rich internal states
without loss) and \emph{epistemic accountability} (how a receiving agent or
orchestrator can formally reason about what it received). Recent work on
latent-space communication (RecursiveMAS, Interlat, KVComm) solves the first
problem but leaves the second entirely open.

We prove that the expressivity--accountability boundary is \textbf{structurally
necessary}: any channel that maximises expressivity in the continuous latent space
is provably non-auditable under standard certification frameworks (DO-178C,
IEC~61508). We then formally characterise the EIP and propose a minimal interface
contract --- the \textbf{CLAIM schema} --- that resolves the problem under bounded
assumptions, using the Transferable Belief Model (TBM, Smets 1994), Belnap's
four-valued logic, and AGM belief revision.

We show that RecursiveMAS's Gradient Death Theorem (Theorem~4.1, Yang et
al.\ 2026) is a special case of a more general result derivable from Information
Bottleneck theory --- \emph{restricted to neural implementations via softmax}, not
to the original IB theory itself --- and that this general result motivates the
need for a formal epistemic interface at the latent/symbolic boundary. A new
Section~5.5 establishes the theoretical unification of the \texttt{belief\_mass}
and \texttt{belnap\_state} fields via Perry \& Tsoukias's (1998) fuzzy Belnap
space, providing complete axiomatic foundations for all five CLAIM invariants.
\end{abstract}

\newpage
\tableofcontents
\newpage

%% ─────────────────────────────────────────────
\section{Introduction}
%% ─────────────────────────────────────────────

\subsection{The Confusion at the Heart of Inter-Agent Communication}

The design of communication protocols between AI agents has, until recently, been
treated as a single problem: how do agents share information efficiently? This
framing masks a fundamental dichotomy.

Consider two agents $A_1$ and $A_2$ collaborating on a classification task.
$A_1$ has processed sensor data and formed an internal representation
$h \in \Real^{d_h}$ encoding its assessment. It must communicate this to $A_2$
and to an orchestrator~$O$ that must synthesise the outputs of multiple agents
into a decision. Two distinct questions arise:

\textbf{Q1 --- Expressivity.} What proportion of $A_1$'s internal state $h$ can be
faithfully transmitted to $A_2$? What is the information loss induced by the
channel?

\textbf{Q2 --- Epistemic accountability.} What does $A_2$ (or $O$) \emph{know}
after receiving $A_1$'s message? Can $O$ formally propagate uncertainty, detect
contradictions between agents, identify the agent responsible for a conflict, and
distinguish ``I observed nothing'' from ``I did not look''?

The existing literature addresses Q1 extensively but Q2 almost entirely in
isolation from Q1. No work formalises the interface between the answer to Q1 (a
latent-space communication mechanism) and the answer to Q2 (a formal epistemic
reasoning framework). This interface --- the \emph{Epistemic Interface} --- is the
subject of this paper.

\subsection{Why This Is Urgent}

Three convergent developments make this formalisation urgent:

\textbf{Development 1 --- Latent communication architectures are maturing.}
RecursiveMAS \cite{yang2026}, Interlat \cite{du2025}, KVComm \cite{ye2025}, and
SDE \cite{tang2025} have demonstrated that agents communicating via latent states
outperform textual MAS by 8--17\% with significant efficiency gains. This
architecture is becoming standard.

\textbf{Development 2 --- Deployment contexts demand auditability.} As MAS expand
into high-stakes domains --- medical diagnosis, autonomous systems, military ISR
(Intelligence, Surveillance, Reconnaissance), financial compliance --- formal
auditability, provenance tracing, and uncertainty quantification become
non-negotiable. These requirements are formalised in DO-178C (avionics software),
IEC~61508 (functional safety), and NATO STANAG protocols.

\textbf{Development 3 --- The gap is not being closed.} Protocols such as A2A,
MCP, and FIPA-ACL standardise message syntax but not epistemic content. Uncertainty
quantification methods for LLMs operate in language space and inherit its
expressive limitations. Deep evidential learning methods \cite{hoarau2025,manchingal2022}
produce formal belief functions but only within individual agents, not across
agent boundaries in a MAS.

\subsection{Contributions}

This paper makes five contributions:

\begin{enumerate}[leftmargin=*]
\item \textbf{Formalisation of the EIP.} We define the EIP precisely as the
problem of constructing a formal mapping $\gamma: \Real^{d_h} \to \Bel(\Theta)$
from agents' latent states to belief structures over a shared frame of discernment
$\Theta$, under auditability, provenance, and qualified silence constraints.

\item \textbf{Generalisation of RecursiveMAS Theorem~4.1.} We show that the
Gradient Death result is a special case of a more general theorem derivable from
Information Bottleneck theory, restricted to neural implementations via softmax ---
not to the IB theory itself.

\item \textbf{Structural necessity of the Expressivity/Accountability boundary.}
We prove that no single channel can simultaneously maximise expressivity (continuous
latent) and satisfy auditability constraints (symbolic, discrete, certifiable).

\item \textbf{The CLAIM schema as a formal solution.} A minimal interface contract
with five invariants satisfying the EIP under bounded assumptions.

\item \textbf{Theoretical unification.} We establish that the \texttt{belief\_mass}
and \texttt{belnap\_state} fields are projections of a single mathematical object
--- a mass function over Perry \& Tsoukias's (1998) fuzzy Belnap space ---
providing complete axiomatic foundations for the CLAIM design choices.
\end{enumerate}

%% ─────────────────────────────────────────────
\section{Background and Related Work}
%% ─────────────────────────────────────────────

\subsection{Latent-Space Communication Between Agents}

\paragraph{Textual MAS.} The dominant paradigm for LLM-based MAS uses natural
language as the communication medium (AutoGen, ChatDev, MetaGPT). This offers
human interpretability but suffers from documented limitations: semantic
compression, information loss at sampling, and inability to transmit complete
reasoning trajectories.

\paragraph{Latent communication.} A growing body of work replaces text with latent
representations: CIPHER transmits probability-weighted token embeddings; AC
transmits last-token hidden states; Interlat extends this to full inter-agent
communication with compression; KVComm shares key-value caches; SDE transmits
state-delta trajectories. RecursiveMAS \cite{yang2026} --- the most systematic ---
treats the entire MAS as a recursive computation in latent space with
co-optimisation via shared gradient credit attribution, demonstrating +8.3\%
accuracy, $1.2\times$--$2.4\times$ speedup, and 34.6\%--75.6\% token reduction.

\paragraph{The theoretical gap.} None of this work addresses what a receiving agent
or orchestrator can formally conclude from the received latent state. The message
is richer --- but the epistemic framework for processing it is no more developed
than for text.

\subsection{Epistemic Frameworks for Individual Agents}

\paragraph{Bayesian approaches.} Standard uncertainty quantification in LLMs uses
probability theory, typically detecting epistemic uncertainty via log-likelihood
thresholding or iterative prompting. These methods remain in language space and
inherit its gradient limitations.

\paragraph{Evidential approaches (Dempster-Shafer / TBM).} DST \cite{shafer1976}
and the TBM \cite{smets1994} represent epistemic uncertainty as mass functions
$m: 2^\Theta \to [0,1]$, enabling explicit representation of ignorance (mass on
$\Theta$) and conflict (mass on $\emptyset$ in the unnormalised formulation). Key
results: Deep EK-NN \cite{hoarau2025} applies $k$-nearest-neighbours in latent
space to produce mass functions by distance-weighted evidential combination,
explicitly separating epistemic uncertainty (non-specificity, Dubois \& Prade)
from aleatoric uncertainty (discord, Klir \& Ramer). Epistemic Deep Learning
\cite{manchingal2022} produces mass functions via random set neural networks.

\paragraph{Critical limitation.} These frameworks produce formal belief functions
\emph{within} individual agents. They do not define how these belief functions
should be communicated across agent boundaries in a MAS, nor how an orchestrator
should combine belief functions from multiple agents while preserving formal
epistemic properties.

\subsection{Communication Protocols for MAS}

FIPA-ACL (2002), A2A (2025), and MCP standardise message envelopes, performatives,
and interaction patterns. They impose no agreement on semantic content or epistemic
representation. The Internet of Agents architecture \cite{ioa2025} proposes
semantic layer L9 for shared ontology negotiation --- an advance toward the EIP,
but it remains agnostic about belief representation, uncertainty, and conflict.
No existing protocol specifies: (a) how an agent expresses its degree of belief in
a proposition, formally and comparably across heterogeneous agents; (b) how a
receiving orchestrator should combine belief expressions from multiple agents;
(c) how to represent qualified silence; (d) how to preserve provenance so that the
origin of a conflicting belief can be identified.

%% ─────────────────────────────────────────────
\section{The Gradient Death Theorem --- Generalisation}
%% ─────────────────────────────────────────────

\subsection{RecursiveMAS Theorem~4.1 --- The Original Result}

\begin{keyresult}
\begin{theorem}[Yang et al.\ 2026 --- Gradient Death in Textual MAS]
Under Assumption~A.1 (textual SFT uses
$R_\text{text}(h) = W_\text{in} \cdot \text{softmax}(W_\text{out} \cdot h)$
as the recursive link, with
$\|W_\text{in}\|_2, \|W_\text{out}\|_2 \le O(1)$),
if tokens are confident with entropy $\varepsilon \le \bar\varepsilon$ where
$\bar\varepsilon \ll 1$, then:
\[
  \left\|\frac{\partial R_\text{text}(h)}{\partial h}\right\|_2 \le O(\varepsilon) \ll 1.
\]
\end{theorem}
\end{keyresult}

\textit{Proof sketch (verified directly in Appendix~A.3 of Yang et al.\ 2026).}
(1) By chain rule: $J_\text{text} = W_\text{in} \cdot S \cdot W_\text{out}$,
where $S = \text{diag}(p) - pp^\top$ and $p = \text{softmax}(W_\text{out} h)$.
(2) By sub-multiplicativity: $\|J_\text{text}\|_2 \le O(\|S\|_2)$.
(3) Since $S$ is the covariance matrix of a categorical distribution (PSD):
$\|S\|_2 \le \text{Tr}(S) = 1 - \|p\|_2^2$.
(4) By the logarithmic inequality $\ln z \le z - 1$:
$\varepsilon \ge \sum_i p_i(-\ln p_i) \ge 1 - \|p\|_2^2$.
(5) Therefore $\|J_\text{text}\|_2 \le O(\varepsilon)$. $\square$

\begin{remark}[Scope correction, v0.2]
The theorem assumes a standard softmax decoder ($T=1$, deterministic). For
temperature $T \ne 1$, the bound becomes $O(\varepsilon/T)$. For stochastic
decoders (top-$k$, top-$p$), the gradient is undefined in the classical sense.
The theorem does not cover stochastic communication channels.
\end{remark}

\subsection{The General Theorem --- Any Discrete Communication Channel}

We prove that gradient death holds for any communication channel that discretises
a continuous latent state.

\begin{definition}[Discrete Communication Channel]
A communication channel $C: \Real^{d_h} \to \Delta(V)$ is a \emph{discrete
communication channel} if it maps a latent state $h$ to a probability distribution
over a finite vocabulary $V$ with $|V| < \infty$. The effective Jacobian of the
round-trip communication cycle (encode $\to$ decode) is
$J_C = \partial R_C(h)/\partial h$,
where $R_C(h) = \mathbb{E}_{v \sim C(h)}[\text{embed}(v)]$ is the expected
re-embedding of the sampled output.
\end{definition}

\begin{keyresult}
\begin{theorem}[Generalised Gradient Death]
Let $C: \Real^{d_h} \to \Delta(V)$ be any discrete communication channel with
output distribution $p_C(h)$. Let $H(p_C(h)) \le \varepsilon$ be the output
entropy of the channel. Under bounded embedding norms, with probability
$\ge 1 - \delta$:
\[
  \left\|\frac{\partial R_C(h)}{\partial h}\right\|_2 \le O(\varepsilon).
\]
\end{theorem}
\end{keyresult}

\textit{Proof sketch.}
The Jacobian decomposes as $J_C = E_\text{embed} \cdot S_C \cdot \nabla_h p_C$,
where $S_C = \text{diag}(p_C) - p_C p_C^\top$ is the covariance matrix of the
output distribution. Since $S_C$ is PSD:
$\|S_C\|_2 \le \text{Tr}(S_C) = 1 - \|p_C\|_2^2$.
By the same logarithmic inequality: $1 - \|p_C\|_2^2 \le H(p_C) \le \varepsilon$.
Under bounded embedding norms: $\|J_C\|_2 \le O(\varepsilon)$. $\square$

\textbf{Key intuition.} The entropy $\varepsilon$ of the channel's output
distribution --- not the vocabulary size, not the encoding mechanism --- determines
the gradient bound. The more confident an agent is ($\varepsilon \to 0$), the more
the gradient vanishes, regardless of the vocabulary, the implementation of the
mapping $h \to V$, or whether the vocabulary represents tokens, ontology concepts,
or JSON schema values.

\subsection{Connection with Information Bottleneck Theory (v0.2 --- Corrected)}

\begin{changebox}
\textbf{Correction v0.1 $\to$ v0.2.}
In v0.1, the connection to IB was presented as if the matrix
$S = \text{diag}(p) - pp^\top$ appeared explicitly in Tishby (1999). This is
incorrect. Tishby's IB optimises directly over the simplex of probability
distributions, without softmax parametrisation. The matrix $S$ is an artefact of
the \emph{neural implementation} of IB --- not of the IB theory itself. This
correction strengthens Theorem~6.1: the Expressivity/Accountability boundary is
architectural, not tied to an implementation choice.
\end{changebox}

In the Information Bottleneck framework \cite{tishby1999}, a communication channel
corresponds to the stochastic mapping $p(T|X)$, where $X$ is the input (latent
state $h$), $T$ is the transmitted representation, and $Y$ is the downstream
prediction target. The IB objective is:
\[
  \min_{p(T|X)} \mathcal{L}_\text{IB} = I(X;T) - \beta \cdot I(T;Y).
\]
The matrix $S = \text{diag}(p) - pp^\top$ does not appear in the original IB
formulation. It appears \emph{only} when the IB channel $p(T|X)$ is parametrised
by a neural network with a softmax layer. This is an artefact of the softmax
parametrisation. RecursiveMAS resolves the problem by replacing the
discrete-softmax channel with a continuous residual channel (RecursiveLink) whose
Jacobian remains close to the identity --- equivalent to solving IB without
saturation, consistent with the GIB principle \cite{westphal2025} of synergistic
feature combination.

\begin{corollary}[Scope of Theorem~4.1]
The Gradient Death result of Yang et al.\ (2026) applies to any discrete channel
--- including JSON-encoded ontologies, FIPA-ACL performatives, and symbolic logic
assertions --- when these channels are implemented via softmax on a categorical
output distribution in a high-confidence regime. It does \emph{not} apply to:
(a) continuous latent channels (RecursiveMAS); (b) stochastic decoders (top-$k$,
top-$p$); (c) algebraic operations on probability masses (TBM combination rules are
not communication channels in the sense of Definition~3.1).
\end{corollary}

\begin{corollary}[Structural Necessity of the Boundary]
No single channel can simultaneously: (a) avoid gradient death (requires continuous
latent space), and (b) satisfy formal auditability constraints (requires symbolic,
discrete, certifiable representations). The boundary between the latent domain and
the epistemic domain is structurally necessary --- not a design choice.
\end{corollary}

%% ─────────────────────────────────────────────
\section{The Epistemic Interface Problem}
%% ─────────────────────────────────────────────

\subsection{Formal Definition}

\begin{definition}[Frame of Discernment]
$\Theta = \{\theta_1, \ldots, \theta_n\}$ is a finite set of mutually exclusive
and exhaustive hypotheses, shared by all agents in the MAS and by the orchestrator.
\end{definition}

\begin{definition}[Belief Mass Function]
A belief mass function $m: 2^\Theta \to [0,1]$ satisfies
$\sum_{A \subseteq \Theta} m(A) = 1$.
In the unnormalised TBM formulation \cite{smets1994}, $m(\emptyset) \ge 0$ is the
conflict mass, preserved rather than redistributed --- avoiding Zadeh's paradox.
\end{definition}

\begin{definition}[Belnap State]
A Belnap state $\beta \in \{T, F, B, N\}$ \cite{belnap1977} represents:
$T$ (True) = proposition believed; $F$ (False) = proposition disavowed;
$B$ (Both) = contradictory evidence received; $N$ (Neither) = no evidence
received --- the qualified-silence state.
\end{definition}

\begin{keyresult}
\begin{definition}[Epistemic Interface Problem (EIP)]
Given agent $A_i$ with internal latent state $h_t \in \Real^{d_h}$ at time $t$,
and an orchestrator $O$ operating over $\Theta$, the EIP is the problem of
constructing a mapping $\gamma_i: \Real^{d_h} \times T \times K \times P \to \Bel(\Theta)$
satisfying simultaneously:
(1) \textbf{Calibration} --- $m$ is a formally calibrated belief function;
(2) \textbf{Qualified silence} --- the system can distinguish ``no evidence''
($\beta = N$) from ``negative evidence'' ($\beta = F$);
(3) \textbf{Provenance} --- the belief function carries a verifiable identifier
of its source agent;
(4) \textbf{Temporal validity} --- $m$ is accompanied by a validity window;
(5) \textbf{Auditability} --- the belief function and its combination rules are
formally certifiable.
\end{definition}
\end{keyresult}

\subsection{Why the EIP Cannot Be Resolved in the Latent Domain}

\begin{proposition}
No continuous latent communication channel simultaneously satisfies constraints
(2)~Qualified Silence, (3)~Provenance, and (5)~Auditability of Definition~4.4.
\end{proposition}

\textit{Argument.} Constraint (5) requires a deterministic, symbolic, certifiable
function --- real-valued vectors are not directly interpretable by certification
frameworks requiring discrete, enumerable, reproducible states. Constraint (2)
requires a ternary distinction impossible to represent via near-zero feature values
without additional symbolic scaffolding. Constraint (3) requires that
\texttt{chain\_id} be explicitly attached --- latent vectors carry no such metadata
in existing architectures. $\square$

\subsection{Why the EIP Cannot Be Resolved by Current Protocols}

\begin{proposition}
No existing agent communication protocol (FIPA-ACL, A2A, MCP, SLIM) resolves the
EIP. In particular: FIPA-ACL includes an \texttt{ontology} field but provides no
mechanism to negotiate or validate belief mass distributions; A2A standardises task
delegation but assumes agents are ``honest about their uncertainty'' without
specifying a valid uncertainty representation; MCP shares context and allows tool
calls but does not specify the epistemic content of tool outputs; none of these
protocols supports $\beta = N$ (qualified silence) as distinct from $\beta = F$
(negative observation); none supports $m(\emptyset)$ as an explicit conflict
indicator; the IoA L9 layer \cite{ioa2025} advances toward semantic negotiation but
remains agnostic about belief representation.
\end{proposition}

%% ─────────────────────────────────────────────
\section{The CLAIM Schema as a Formal Solution}
%% ─────────────────────────────────────────────

\subsection{Definition}

\begin{keyresult}
\begin{definition}[CLAIM --- Contextual Latent Attestation of Indexed Meaning]
A CLAIM is a 6-tuple:
\begin{itemize}[noitemsep]
\item \texttt{proposition}: $\Phi$ --- propositional content;
\item \texttt{belief\_mass}: $m: 2^\Theta \to [0,1]$ --- unnormalised TBM mass
function, $m(\emptyset)$ = explicit conflict mass;
\item \texttt{belnap\_state}: $\beta \in \{T, F, B, N\}$ --- four-valued epistemic
state;
\item \texttt{illocution}: $\kappa \in \{\text{OBSERVE, INFER, DEDUCE, ASSUME}\}$
--- speech act type;
\item \texttt{freshness}: $(t_\text{obs}, \Delta t_\text{valid})$ --- timestamp +
validity window;
\item \texttt{provenance}: \texttt{chain\_id} --- PROV-O compatible identifier.
\end{itemize}
\end{definition}
\end{keyresult}

\subsection{EIP Satisfaction}

\begin{theorem}[CLAIM Satisfies the EIP]
Under Assumption~5.1 (the mapping $\gamma_i: \Real^{d_h} \to \text{CLAIM}$ is
implemented by a calibrated, deterministic function for each agent $A_i$), the
CLAIM schema satisfies the five constraints of Definition~4.4.
\end{theorem}

\textit{Proof.}
(1) Calibration: via TBM mass function on annotated $\Theta$ data.
(2) Qualified silence: $\beta = N$, explicitly distinct from $B$, $T$, $F$.
(3) Provenance: \texttt{chain\_id} implementing PROV-O lineage.
(4) Temporal validity: DST discounting $m_\text{disc}(A) = \alpha \cdot m(A) + (1-\alpha) \cdot m(\Theta)$.
(5) Auditability: the CLAIM is a stateless, deterministic data structure whose
combination rules are algebraic operations with formal proofs. $\square$

\subsection{Orchestrator Rules --- Axiomatic Justification (v0.2)}

\paragraph{Rule O1 --- TBM Conjunctive Combination \cite{smets1994}.}
\[
  m_{1 \oplus 2}(A) = \sum_{B \cap C = A} m_1(B) \cdot m_2(C), \quad A \subseteq \Theta.
\]
\textbf{Axiomatic justification} (Klawonn \& Smets 1992; Pichon \& Deno\oe{}ux 2008):
the unnormalised conjunctive rule is the \emph{unique} rule satisfying simultaneously
(a) associativity, (b) commutativity, (c) neutral element (the vacuous belief
function), (d) the Least Commitment Principle (LCP) --- it never adds more belief
than the available information justifies. Dempster's normalised rule satisfies
(a)--(c) but violates LCP by redistributing $m(\emptyset)$ (Zadeh's paradox).
Rule O1 is therefore axiomatically necessary, not arbitrary. The conflict mass
$m(\emptyset)$ is preserved as an explicit indicator of inter-agent disagreement.

\paragraph{Rule O3 --- Source Deduplication (Cautious Rule, Deno\oe{}ux 2008).}
Activated when common \texttt{chain\_id} prefixes detect a shared source between agents.
The conjunctive cautious rule is idempotent: $m \oplus_\text{cautious} m = m$. This
property is necessary when agents share an undeclared common source: combining the
same evidence twice must not double-count it. The combination O1 (independence) +
O3 (dependence) provides complete coverage of the combination problem. To our
knowledge, this Smets/Deno\oe{}ux pair has not been explicitly formulated as an
adaptive combination policy based on detected provenance in the existing literature.

\paragraph{Rule O5 --- Human Escalation Threshold.}
When $m(\emptyset) \ge \theta_\text{conflict}$, the orchestrator suspends automated
decision-making and escalates to a human operator rather than resolving the conflict
algorithmically. This is not a conservative choice --- it is a consequence of the
LCP. Redistributing $m(\emptyset)$ would violate LCP (adding unjustified belief).
The only decision compatible with LCP in the face of persistent conflict is to
expose it to the decision-maker. Escalation transmits: the value of $m(\emptyset)$,
the source pair identified by PCR5, Belnap states of all involved agents, and
illocution types (to distinguish conflicts between observations, inferences, or
assumptions).

\subsection{The Calibration Problem --- The Open Question}

Theorem~5.1 proves that CLAIM satisfies the EIP under Assumption~5.1 --- that the
mapping $\gamma_i: \Real^{d_h} \to \text{CLAIM}$ is calibrated and deterministic.
This assumption is non-trivial and represents the main open problem identified by
this paper.

Deep EK-NN \cite{hoarau2025} provides the most complete answer for classification
tasks: use $k$-nearest-neighbours in latent space to construct mass functions by
distance-weighted evidential combination. This approach: (\checkmark) produces
calibrated mass functions from latent representations; (\checkmark) explicitly
separates epistemic from aleatoric uncertainty; (\checkmark) works with pretrained
models; ($\times$) requires a labelled reference set in latent space; ($\times$)
does not extend to open-vocabulary generation tasks; ($\times$) does not address
the inter-agent boundary.

\subsection{Theoretical Unification via the Fuzzy Belnap Space (v0.2 --- New)}

\subsubsection{The Fuzzy Belnap Space (Perry \& Tsoukias 1998)}

Perry \& Tsoukias \cite{perry1998} define a continuous extension of Belnap's
four-valued logic. The truth-value space becomes the set of convex combinations of
the four reference values:
\[
  \mathbb{M} = \{(t, f, k, u) \in [0,1]^4 \mid t + f + k + u = 1\}
\]
where $t$ = degree of truth, $f$ = degree of falsity, $k$ = degree of contradiction
(both), $u$ = degree of unknown (neither). The four classical Belnap values are the
corners of the simplex: $T = (1,0,0,0)$, $F = (0,1,0,0)$, $B = (0,0,1,0)$,
$N = (0,0,0,1)$.

\subsubsection{Isomorphism with TBM Mass Functions}

A mass function $m$ over $\Theta_\text{Belnap} = \{T, F, B, N\}$ assigns masses
to the 16 subsets of $\{T, F, B, N\}$. The atomic masses ---
$m(\{T\}), m(\{F\}), m(\{B\}), m(\{N\})$ --- are exactly the components
$(t, f, k, u)$ of the Perry-Tsoukias space when restricted to masses on singletons.
This is not an analogy --- it is an isomorphism: every point of $\mathbb{M}$
corresponds to a belief mass function concentrated on the four Belnap atoms.
The unified object is:
\[
  m_\text{meta}: 2^{\Theta \times \{T,F,B,N\}} \to [0,1]
\]
This object simultaneously encodes: (a) what the agent believes about the world
(hypotheses $\theta \in \Theta$), and (b) the agent's epistemic status regarding
its own belief (truth, falsity, contradiction, ignorance).

\subsubsection{Projections and CLAIM Fields}

The two CLAIM fields are projections of $m_\text{meta}$:
\begin{itemize}[noitemsep]
\item \texttt{belief\_mass} (world level): $m_\text{world}(\theta) = \sum_\beta m_\text{meta}(\theta, \beta)$ --- marginalisation over epistemic status.
\item \texttt{belnap\_state} (epistemic level): $m_\text{epist}(\beta) = \sum_\theta m_\text{meta}(\theta, \beta)$ --- marginalisation over world hypotheses.
\end{itemize}
The discrete $\beta \in \{T, F, B, N\}$ is the limit case where $m_\text{epist}$
is concentrated on a single atom. The full $m_\text{meta}$ object is an
architectural option for NMT 4--5.

\subsubsection{Connection with O-Information}

A remarkable connection with the synergy measure (GIB, \cite{westphal2025}) emerges:

\begin{center}
\begin{tabular}{lll}
\toprule
O-Information & Belnap State & TBM Interpretation \\
\midrule
$+1$ (max synergy) & $T$ (true) & Pure complementarity --- information combines perfectly \\
$-1$ (max redundancy) & $F$ (false) & Total overlap --- agents duplicate \\
$0$ with mixed components & $B$ (contradiction) & Tension --- synergy and redundancy coexist \\
$0$ with null components & $N$ (ignorance) & Independence --- no shared information \\
\bottomrule
\end{tabular}
\end{center}

\subsubsection{Adaptive Stopping Criterion for RecursiveMAS}

The fuzzy Belnap representation immediately suggests a principled stopping
criterion for RecursiveMAS recursive rounds:

\textbf{Stopping criterion:} Terminate recursion when $u$ (unknown) and $k$
(contradiction) fall below their respective thresholds $\theta_u$ and $\theta_k$
--- indicating that belief has stabilised and additional rounds will no longer
reduce epistemic uncertainty.

\textbf{Escalation criterion:} If $k$ (contradiction) remains persistently above
$\theta_k$ despite $r \ge r_\text{max}$ rounds, trigger Rule O5 (human escalation)
--- the system has identified an irresolvable inter-agent conflict.

This criterion is formally justified by the axiomatic structure of the CLAIM:
stopping when $u \downarrow$ and $k \downarrow$ is a direct application of the LCP
(believe only what the information justifies; stop when no new information arrives).
To our knowledge, no existing MAS paper proposes a formally justified stopping
criterion for recursive collaboration based on epistemic state.

%% ─────────────────────────────────────────────
\section{The Expressivity / Accountability Boundary}
%% ─────────────────────────────────────────────

\begin{keyresult}
\begin{theorem}[Structural Necessity of the Expressivity/Accountability Boundary]
No communication channel $C: \Real^{d_h} \to O$ simultaneously satisfies:
(a) \textbf{Gradient preservation}: $\exists\, c > 0$, $\forall h \in M$,
$\|\partial R_C / \partial h\|_2 \ge c$ for all entropy values $\varepsilon$; and
(b) \textbf{Formal auditability}: $O$ is a discrete, enumerable, certifiable output
space.
\end{theorem}
\end{keyresult}

\textit{Proof.}
By Theorem~3.1, if $O$ is discrete (condition b), then
$\|\partial R_C / \partial h\|_2 \le O(\varepsilon)$, which vanishes as
$\varepsilon \to 0$ (violates condition a). If $O$ is continuous (satisfies
condition a via RecursiveLink), then $O = \Real^{d'}$ is not a discrete,
enumerable space (violates condition b). $\square$

\begin{center}
\begin{tabular}{p{5.5cm}p{5.5cm}}
\toprule
\textbf{Continuous Latent Domain} & \textbf{Symbolic Epistemic Domain} \\
\midrule
RecursiveMAS, Interlat, KVComm & CLAIM + Orchestrator \\
$h_t \in \Real^{d_h}$ & $m: 2^\Theta \to [0,1]$, $\beta \in \{T,F,B,N\}$ \\
Gradient preservation & Certifiable, auditable \\
Expressivity maximisation & Accountability maximisation \\
Opaque to certification & DO-178C / IEC 61508 compatible \\
\midrule
\multicolumn{2}{c}{$\longleftarrow \gamma_i \longrightarrow$ (open calibration problem --- Section~5.4)} \\
\bottomrule
\end{tabular}
\end{center}

%% ─────────────────────────────────────────────
\section{Related Work --- Positioning}
%% ─────────────────────────────────────────────

\begin{center}
\begin{tabular}{lccc}
\toprule
Work & Expressivity & EIP & Across Agent Boundaries \\
\midrule
RecursiveMAS \cite{yang2026} & \checkmark Complete & $\times$ & N/A (latent domain) \\
Interlat \cite{du2025} & \checkmark Partial & $\times$ & N/A \\
Deep EK-NN \cite{hoarau2025} & N/A & \checkmark Single agent & $\times$ \\
Epistemic DL \cite{manchingal2022} & N/A & \checkmark Single agent & $\times$ \\
FIPA-ACL / A2A / MCP & $\times$ (syntax) & $\times$ & Syntactic only \\
IoA L9 \cite{ioa2025} & N/A & $\sim$ Partial (ontology) & $\checkmark$ Non-epistemic \\
\textbf{This work --- CLAIM} & N/A (by Thm.~6.1) & \checkmark Formal & \checkmark Defined \\
\bottomrule
\end{tabular}
\end{center}

%% ─────────────────────────────────────────────
\section{Conclusion}
%% ─────────────────────────────────────────────

We have identified, formalised, and partially resolved the Epistemic Interface
Problem in multi-agent systems. The central contributions of v0.2 are:

\begin{enumerate}[leftmargin=*]
\item The EIP is a problem distinct from expressivity. Improving the communication
channel (RecursiveMAS) does not resolve the EIP. Improving epistemic frameworks
within individual agents (Deep EK-NN) does not resolve the EIP. The interface
between the two domains is the gap.

\item Gradient death is more general than RecursiveMAS claims --- it applies to
any discrete channel parametrised by softmax --- but is not predicted by the IB
theory itself. This correction strengthens the paper's argument by making it
implementation-agnostic at the architectural level.

\item The Expressivity/Accountability boundary is necessary (Theorem~6.1).
Hybrid neuro-symbolic architectures must make this boundary explicit.

\item CLAIM is a minimal formal solution satisfying the EIP under bounded
assumptions. The calibration problem ($\gamma_i$) is the open research frontier.

\item The \texttt{belief\_mass} and \texttt{belnap\_state} fields are theoretically
unified: they are projections of a single mass function over Perry \& Tsoukias's
(1998) fuzzy Belnap space. This provides complete axiomatic foundations for all
CLAIM invariants and for orchestrator rules O1/O3/O5.
\end{enumerate}

%% ─────────────────────────────────────────────
\section{Open Problems}
%% ─────────────────────────────────────────────

\begin{itemize}[leftmargin=*]
\item Formal convergence proof for the proposed $\gamma_i$ calibration (Section~5.4)
\item Extension to open-world hypothesis spaces ($|\Theta|$ not predefined)
\item Empirical validation on multimodal ISR datasets
\item Integration with formal verification tools for orchestrator combination rules
\item Strategic behaviour analysis: what happens if agents have incentives to
misreport their \texttt{belief\_mass}?
\item Formal study of the Smets/Deno\oe{}ux duality (O1/O3) as an adaptive
combination policy
\end{itemize}

%% ─────────────────────────────────────────────
\section{Future Extensions}
%% ─────────────────────────────────────────────

Three concrete research directions follow from the theoretical unification of
Section~5.5.

\subsection{Latent Channel with Credal Output}

Replace RecursiveMAS's final text output with a CLAIM: the last RecursiveLink
produces not a decoded text but a quadruplet $(t, f, k, u)$ in the fuzzy Belnap
space --- a continuous credal output. Residual epistemic uncertainty is measured by
$u$ (unknown) and $k$ (contradiction). This architecture is the direct
implementation of Corollary~3.2: the continuous latent domain (RecursiveMAS) feeds
the symbolic epistemic domain (CLAIM) via the calibrated mapping $\gamma_i$.

\subsection{Recursive Belief Fusion}

After each RecursiveMAS round, combine all agents' CLAIM mass functions by an
adaptive rule: Rule O1 (Smets, unnormalised conjunctive) for independent agents
(distinct \texttt{chain\_id} prefixes); Rule O3 (Deno\oe{}ux cautious) for
source-correlated agents (common \texttt{chain\_id} prefixes). This adaptive fusion
transforms RecursiveMAS from a purely latent architecture into an explicit epistemic
reasoning system. The accumulated conflict mass $m(\emptyset)$ across rounds is the
signal for both additional rounds (if below $\theta_\text{conflict}$) and human
escalation (if above $\theta_\text{conflict}$).

\subsection{Adaptive Stopping via Epistemic Uncertainty Measure}

Using the quadruplet $(t, f, k, u)$ as the agent state, define a formally justified
stopping criterion:
\begin{itemize}[noitemsep]
\item Continue recursion when $u + k > \theta_\text{uncertainty}$
\item Stop when $u + k \le \theta_\text{uncertainty}$ for all agents (belief stabilised)
\item Escalate when $k > \theta_\text{conflict}$ for $r \ge r_\text{max}$ consecutive rounds
\end{itemize}
This is the first formally justified adaptive stopping rule for recursive
multi-agent collaboration, grounded in TBM axiomatics and the LCP. It replaces
RecursiveMAS's fixed-round approach ($r = 1, 2, 3$) with a principled termination
condition.

\textbf{Integration path:} NMT 2--3: operational CLAIM schema with discrete Belnap
(current specification). NMT 4--5: full $m_\text{meta}: 2^{\Theta \times
\{T,F,B,N\}} \to [0,1]$ with continuous fuzzy Belnap output, adaptive stopping and
recursive belief fusion.

%% ─────────────────────────────────────────────
\bibliographystyle{plainnat}
\bibliography{references}

\end{document}
